% Embedded version: Nordic, NXP, Arm, TomTom, Bosch, Continental,
% Valeo and other automotive / firmware employers.
% Build: pdflatex cv-embedded.tex (twice). Needs the lato package.
\documentclass[10pt]{article}
\usepackage{cv-style}
\hypersetup{pdftitle={Abdelrahman Khalifa - Embedded Software Engineer - CV},pdfauthor={Abdelrahman Khalifa}}

\begin{document}

\cvheader{Abdelrahman Khalifa}{Embedded Software Engineer --- Zephyr $\cdot$ Embedded Linux $\cdot$ Rust}

\section*{Summary}
Embedded software engineer with over a year of professional firmware work on Zephyr RTOS and FreeRTOS, embedded Linux experience with Yocto and Buildroot, and upstream contributions to Zephyr. Writes C, modern C++ and Rust; completed Google Summer of Code 2026 with The Linux Foundation.

\section*{Experience}
\cventry{Embedded Systems Engineer --- EKSON}{Aug 2025 -- Present}
\begin{itemize}
  \item Developed Zephyr RTOS firmware for an ESP32-S3 3-DoF motion simulator, migrating it from FreeRTOS for more deterministic behavior and production readiness.
  \item Integrated an MPU6050 IMU over I2C with calibration, filtering and compensation for stable motion estimation.
  \item Built a modern C++/Qt desktop control application with reliable UART control of the simulator and optimized its performance by \textbf{\textasciitilde40\%} through asynchronous processing and fewer threads.
\end{itemize}
\cventry{Google Summer of Code 2026 Contributor --- The Linux Foundation / OpenPrinting}{May 2026 -- Oct 2026}
\cvsub{COSMIC desktop printer setup tool in Rust \enspace|\enspace \cvlink{https://github.com/Abd002/GSOC-2026}{github.com/Abd002/GSOC-2026} \enspace|\enspace \cvlink{https://github.com/Abd002/cosmic-printers}{github.com/Abd002/cosmic-printers}}
\begin{itemize}
  \item Built a Linux printer-management stack as 5 Rust crates with an async client, a CUPS/IPP/DNS-SD backend behind a varlink IPC service, and a shared UI; tested on real HP printers.
  \item Wrote FFI bindings from Rust to C (libcups3 DNS-SD API) and fixed crashes in libcups' Avahi code; \textbf{7 PRs merged upstream}.
\end{itemize}
\cventry{Embedded Linux Systems Intern --- ODC}{Dec 2023 -- Jan 2024}
\begin{itemize}
  \item Built Linux systems with Yocto and Buildroot; worked with kernel modules, device drivers, device trees, file systems and hardware flashing.
\end{itemize}

\section*{Open Source}
\begin{itemize}
  \item \textbf{Zephyr RTOS} (Mar 2024): implemented the \texttt{IP\_MULTICAST\_LOOP} socket option for IPv4 across 6 networking modules (150+ lines), reviewed and merged upstream (\cvlink{https://github.com/zephyrproject-rtos/zephyr/commit/b11703623c02f5afa4a5b0b7abdce3e324d7d470}{commit b117036}).
  \item \textbf{OpenPrinting libcups and cups-rs} (2026): fixed DNS-SD crashes in libcups (C); ported cups-rs to libcups3 with CI (\cvlink{https://github.com/OpenPrinting/cups-rs/pull/18}{PRs \#18--\#22}); Rust \texttt{lp}/\texttt{lpstat} in Winter of Code 2026.
  \item \textbf{Rust ecosystem} (Dec 2025): Boa JavaScript engine string optimization (\cvlink{https://github.com/boa-dev/boa/pull/4553}{PR \#4553}) and a Rust Clippy lint fix (\cvlink{https://github.com/rust-lang/rust-clippy/pull/16214}{PR \#16214}).
\end{itemize}

\section*{Projects}
\begin{itemize}
  \item \textbf{AUTOSAR-Based V2X Ecosystem for SDV} (graduation project; FreeRTOS, AUTOSAR): software-defined vehicle prototype with ADAS (LKAS, driver monitoring), FOTA updates, and a secure bootloader with rollback. \cvlink{https://github.com/Abd002/Graduation-Project}{GitHub}
  \item \textbf{Seat Heater Control System} (Tiva C, FreeRTOS): real-time controller with temperature monitoring, heater-level control and UART diagnostics, using tasks, mutexes and event flags. \cvlink{https://github.com/Abd002/SeatHeaterControlSystem}{GitHub}
  \item \textbf{IoT Device Communication System} (C++, Raspberry Pi 4, QEMU, Yocto): TCP unicast and UDP multicast framework between a Pi server and emulated clients; Yocto cross-compilation and Make. \cvlink{https://github.com/Abd002/IoT-Comm-System-RPi4-QEMU}{GitHub}
  \item \textbf{AUTOSAR PORT and GPIO drivers} (C, TM4C): developed and validated MCAL drivers following the AUTOSAR specification.
  \item \textbf{Traffic Light Controller} (Verilog, Spartan-3E FPGA): FSM with debounce and delay modules, cutting signal timing errors by 20\%. \cvlink{https://github.com/Abd002/Traffic-Light-Controller}{GitHub}
  \item \textbf{Door Locker Security System} (C, two MCUs): password lock with keypad, LCD, I2C EEPROM, timer and DC motor; UART between the MCUs. \cvlink{https://github.com/Abd002/Door-Locker-Security-Systems}{GitHub}
\end{itemize}

\section*{Technical Skills}
\cvskill{Languages}{C, C++ (modern C++, Qt), Rust, Python, Bash, Verilog/VHDL}
\cvskill{RTOS \& firmware}{Zephyr, FreeRTOS, ARM Cortex-M (STM32, TM4C), ESP32, AVR, AUTOSAR, MISRA C, bootloaders, FOTA}
\cvskill{Embedded Linux}{Linux (Ubuntu/Debian), operating systems, Yocto, Buildroot, kernel modules, device drivers, device trees, QEMU, Raspberry Pi}
\cvskill{Protocols \& tools}{UART, SPI, I2C, CAN, LIN, USB, Ethernet, TCP/UDP; Git, CMake, Make, CI, testing, debugging, performance optimization, code review}

\section*{Education}
\cventry{B.Sc. Computer and Systems Engineering --- Minia University}{2020 -- 2025}
Very Good (84.6\%) with Honors. Codeforces Expert and 4-time ECPC finalist; Deloitte Software Engineering Mentorship (2024).

\end{document}
